\documentclass[11pt]{article}

\usepackage[margin=1.05in]{geometry}
\usepackage{amsmath,amssymb,amsthm,mathtools}
\usepackage{enumitem}
\usepackage[colorlinks=true,linkcolor=blue,citecolor=blue,urlcolor=blue]{hyperref}

\title{Residual-Mass Transform Order for LRU Caching\ and the Coupon Collector}
\author{Christopher D. Long}
\date{June 2026}

\newtheorem{theorem}{Theorem}[section]
\newtheorem{lemma}[theorem]{Lemma}
\newtheorem{proposition}[theorem]{Proposition}
\newtheorem{corollary}[theorem]{Corollary}
\newtheorem{definition}[theorem]{Definition}
\newtheorem{remark}[theorem]{Remark}
\newtheorem{example}[theorem]{Example}

\newcommand{\E}{\mathbb E}
\newcommand{\Pbb}{\mathbb P}
\newcommand{\R}{\mathbb R}
\newcommand{\N}{\mathbb N}
\newcommand{\dd}{\,d}
\newcommand{\1}{\mathbf 1}
\newcommand{\DeltaN}{\Delta_N}
\newcommand{\Exp}{\operatorname{Exp}}
\newcommand{\MR}{\operatorname{MR}}
\newcommand{\Hit}{\operatorname{Hit}}
\newcommand{\Var}{\operatorname{Var}}

\allowdisplaybreaks

\begin{document}
\maketitle

\begin{abstract}
For the independent reference model with a positive popularity vector
$p\in\Delta_N^\circ$, let $S_C(p)$ be the first $C$ distinct items in size-biased order,
equivalently the stationary LRU cache set of capacity $C$.  Put
\[
        M_C(p)=p(S_C(p)),\qquad Q_C(p)=1-M_C(p).
\]
Thus $M_C$ is the random cached mass and $Q_C$ is the random residual miss mass.
We prove a finite-dimensional radial transform-order theorem.  Along every nonconstant
ray $p_\theta=u+\theta(p-u)$ from the uniform law, $M_C(p_\theta)$ is increasing in
increasing-convex order.  Equivalently, for every decreasing convex cost $\psi$,
\[
        \theta\mapsto \E\psi(Q_C(p_\theta))
\]
is nondecreasing, and is strict for the natural strictly monotone costs.  The proof is a
pathwise crossing argument using a common Gumbel representation of the Plackett--Luce
size-biased order: higher-than-uniform coordinates can only overtake lower-than-uniform
coordinates as $\theta$ increases.

The result simultaneously strengthens radial LRU extremality and gives a new exact
coupon-collector transfer principle.  In LRU, it implies monotonicity not only of the
mean hit and miss rates but also of no-miss survival probabilities.  In coupon collecting,
it implies that every one-step discovery increment $T_{C+1}-T_C$ is radially increasing
in stochastic order, and that
\[
        \theta\mapsto \E Q_C(p_\theta)^{-s}
\]
is strictly increasing for all $s>0$.  The exact residual-mass identities
\[
        \MR_{LRU}(C;p)=\E Q_C(p),
        \qquad
        \E_p[T_{C+1}-T_C]=\E Q_C(p)^{-1}
\]
therefore become two projections of a single ordered residual-hazard law.  This clarifies
and strengthens the asymptotic LRU--coupon-collector relation of Berthet.
\end{abstract}

\tableofcontents

\section{Introduction}

The independent reference model (IRM) is the standard model in which successive requests
are independent with common law
\[
        p=(p_1,\ldots,p_N)\in\Delta_N^\circ,
        \qquad
        p_i>0,
        \qquad
        \sum_i p_i=1.
\]
The same popularity vector governs two classical objects.

First, for an LRU cache of capacity $C$, the stationary cache set is the set of the first
$C$ distinct items encountered when scanning the IRM sequence backwards from stationarity.
Equivalently, it is the set of the first $C$ items in a size-biased permutation with
weights $p$.  Second, in the coupon collector, the set of the first $C$ distinct coupons
seen in a forward run has the same size-biased distribution.  Thus both models share the
same random prefix set $S_C$.

The central random variable of this paper is the residual mass
\[
        Q_C=1-p(S_C).
\]
Its complement
\[
        M_C=p(S_C)
\]
is the random mass already represented in the size-biased prefix.  In LRU language,
$M_C$ is the conditional hit probability given the current cache set and $Q_C$ is the
conditional miss probability.  In coupon-collector language, $Q_C$ is the conditional
success probability for discovering a new coupon after $C$ distinct coupons have already
been collected.

The exact identities
\begin{equation}\label{eq:basic-identities-intro}
        \MR_{LRU}(C;p)=\E Q_C(p),
        \qquad
        \E_p[T_{C+1}-T_C]=\E Q_C(p)^{-1}
\end{equation}
are elementary, but they are structurally revealing.  They show that the LRU miss rate
and the coupon-collector waiting increment are the positive and reciprocal first moments
of the same random residual mass.  The product relation
\[
        \MR_{LRU}(C;p)\,\E_p[T_{C+1}-T_C]
        =\E Q_C\,\E Q_C^{-1}\ge 1
\]
is therefore a Jensen gap.  Equality holds at uniform popularity; away from uniform it
measures dispersion of the residual mass.

Berthet \cite{Berthet2017} derived and emphasized an asymptotic relation of the form
\[
        \MR_{LRU}(C;p)\,\bigl(\E T_{C+1}-\E T_C\bigr)\approx 1
\]
between LRU miss rates and coupon-collector waiting-time increments.  The identities
above show that the approximation is best understood as a concentration statement for
$Q_C$, not as an exact reciprocal relation between the two expectations.  The present
paper strengthens this viewpoint by proving a transform-order theorem for $Q_C$ along
rays from uniform.

Let
\[
        u=(1/N,\ldots,1/N),\qquad
        p_\theta=u+\theta(p-u),\qquad 0\le \theta\le 1.
\]
Our main result says that, for every $1\le C\le N-1$, the cached mass
$M_C(p_\theta)$ is increasing in increasing-convex order.  Equivalently, for every
decreasing convex function $\psi$,
\[
        \theta\mapsto \E\psi(Q_C(p_\theta))
\]
is nondecreasing.  This single statement gives, among other consequences,
\[
        \E Q_C(p_\theta)\downarrow,
        \qquad
        \E e^{-\lambda Q_C(p_\theta)}\uparrow,
        \qquad
        \E Q_C(p_\theta)^{-s}\uparrow,
        \qquad
        \E Q_C(p_\theta)^r\downarrow\quad(0<r\le1).
\]
Thus the LRU miss cost $q$ and the coupon-collector reciprocal cost $q^{-1}$ are two
special costs in a larger ordered family.

\subsection*{Relation to earlier LRU radial extremality}

A previous paper \cite{LongLRU2026} proved that the exact stationary LRU hit rate
$\Hit_C(p)=\E M_C(p)$ is strictly increasing along every nonconstant ray from uniform.
That result used the exponential-age representation of LRU and produced an explicit
positive pair-square radial derivative formula.  It also established the radial
restriction of the Fill--Holst conjecture for move-to-front search-cost tails and
explained why full Schur-convexity fails.

The present theorem strictly extends the conclusion of \cite{LongLRU2026}: the hit-rate
monotonicity is the linear case $\phi(x)=x$ of increasing-convex order for $M_C$.
However, the earlier positive-kernel proof still has independent value.  It gives an
explicit local derivative certificate and computable pair kernels; it also isolates the
finite-dimensional algebraic positivity that survives on rays even though arbitrary
majorization directions can fail.  The Gumbel crossing method developed here gives a
shorter and stronger global order theorem, while the pair-kernel proof remains useful for
sensitivity analysis, local certificates, and transfer to models where a common-Gumbel
ranking representation is unavailable.

\subsection*{Related work}

Flajolet, Gardy and Thimonier \cite{FlajoletGardyThimonier1992} gave a unified analytic
framework for the birthday paradox, coupon collectors, caching algorithms and
self-organizing search under general nonuniform distributions.  Exact and asymptotic LRU
analysis under the IRM goes back at least to King, Fagin and Fagin--Price
\cite{King1971,Fagin1977,FaginPrice1978}, with later asymptotic work on move-to-front
and LRU fault probabilities by Jelenkovic \cite{Jelenkovic1999}.  Che, Tung and Wang
\cite{CheTungWang2002} introduced the characteristic-time approximation that is now a
standard engineering approximation for LRU.  Boneh and Hofri \cite{BonehHofri1997}
surveyed coupon-collector formulae and computational methods.  Anceaume, Busnel,
Schulte-Geers and Sericola \cite{AnceaumeBusnelSericola2016} proved stochastic
optimization results for generalized coupon collection, including almost-uniform
minimization for the time to collect a fixed number of distinct coupons.

The present contribution is different in focus: it identifies an exact finite-$N$
residual-hazard law common to LRU and the coupon collector, and proves a radial transform
order for that law.

\section{Size-biased order, LRU, and coupon collection}

\subsection{Size-biased permutations}

Let $p\in\Delta_N^\circ$.  A size-biased permutation with weights $p$ is the random
permutation $(\Pi_1,\ldots,\Pi_N)$ satisfying
\begin{equation}\label{eq:pl-law}
        \Pbb(\Pi_1=i_1,\ldots,\Pi_N=i_N)
        =\prod_{r=1}^N \frac{p_{i_r}}{1-\sum_{s<r}p_{i_s}},
\end{equation}
for every ordering $(i_1,\ldots,i_N)$ of $[N]$.  This is the
Plackett--Luce distribution.

For $0\le C\le N$, put
\[
        S_C=\{\Pi_1,\ldots,\Pi_C\},
\]
with $S_0=\varnothing$ and $S_N=[N]$.  Define
\[
        M_C=p(S_C)=\sum_{i\in S_C}p_i,
        \qquad
        Q_C=1-M_C.
\]
The nontrivial capacities are $1\le C\le N-1$.  At $C=0$, $Q_0=1$; at $C=N$,
$Q_N=0$.

\begin{lemma}[Exponential and Gumbel representations]\label{lem:gumbel}
Let $E_1,\ldots,E_N$ be independent exponential random variables with rate one.  Then the
increasing order of
\[
        A_i=\frac{E_i}{p_i}
\]
has the size-biased law \eqref{eq:pl-law}.  Equivalently, if
\[
        G_i=-\log E_i,
\]
then $G_i$ are independent standard Gumbel variables, and the decreasing order of
\[
        Y_i=G_i+\log p_i
\]
has the size-biased law.
\end{lemma}

\begin{proof}
The variables $A_i$ are independent exponentials with rates $p_i$.  The first clock to
ring has label $i$ with probability $p_i/\sum_jp_j=p_i$.  Conditional on the first label,
the residual clocks of the unchosen items are independent exponentials with their original
rates, by memorylessness.  Iterating gives \eqref{eq:pl-law}.  Since increasing order of
$E_i/p_i$ is the same as decreasing order of $-\log E_i+\log p_i$, the Gumbel form follows.
\end{proof}

\subsection{LRU cache sets}

Under the IRM, scan the stationary request sequence backwards from a reference time.  The
first distinct item encountered is distributed as $p$.  Given that a set $A$ has already
been encountered, the next new item is $i\notin A$ with probability
\[
        \frac{p_i}{1-p(A)}.
\]
Therefore the backward order of distinct items is a size-biased permutation.  The LRU
cache of capacity $C$ contains exactly the first $C$ distinct items in this backward scan.
Thus its stationary cache set has the same law as $S_C$.

Consequently,
\begin{equation}\label{eq:lru-hit-miss}
        \Hit_{LRU}(C;p)=\E M_C(p),
        \qquad
        \MR_{LRU}(C;p)=\E Q_C(p).
\end{equation}
Moreover, conditional on the stationary cache set, the next request hits with probability
$M_C$ and misses with probability $Q_C$.  Hence, in discrete time,
\begin{equation}\label{eq:no-miss-discrete}
        \Pbb(\hbox{the next } n \hbox{ requests are all LRU hits})
        =\E M_C^n,
        \qquad n\ge0.
\end{equation}
In the rate-one Poissonized IRM,
\begin{equation}\label{eq:no-miss-continuous}
        \Pbb(\hbox{no LRU miss before time }t)=\E e^{-tQ_C}.
\end{equation}

\subsection{Coupon-collector increments}

Let $T_C$ be the number of draws needed to have seen $C$ distinct coupon types, with
$T_0=0$.  The set of the first $C$ collected types is distributed as $S_C$.  Conditional on
$S_C$, the probability that the next draw is a new coupon is $Q_C$.  Therefore
\[
        Y_C:=T_{C+1}-T_C
\]
is conditionally geometric on $\{1,2,\ldots\}$ with success probability $Q_C$:
\begin{equation}\label{eq:geom-survival}
        \Pbb(Y_C>n\mid S_C)=(1-Q_C)^n=M_C^n,
        \qquad n\ge0.
\end{equation}
Thus
\begin{equation}\label{eq:increment-mean}
        \E[Y_C]=\E Q_C^{-1}.
\end{equation}
In the Poissonized coupon collector, the corresponding waiting time from $C$ to $C+1$
distinct coupons is conditionally exponential with rate $Q_C$, and therefore has survival
\begin{equation}\label{eq:poissonized-increment-survival}
        \Pbb(Y_C^{\rm cont}>t)=\E e^{-tQ_C}.
\end{equation}

\begin{proposition}[Exact residual-mass identities]\label{prop:residual-identities}
For every $p\in\Delta_N^\circ$ and $0\le C<N$,
\begin{equation}\label{eq:residual-identities}
        \MR_{LRU}(C;p)=\E Q_C(p),
        \qquad
        \E_p[T_{C+1}-T_C]=\E Q_C(p)^{-1}.
\end{equation}
Consequently, for $1\le C<N$,
\begin{equation}\label{eq:jensen-product}
        \MR_{LRU}(C;p)\,\E_p[T_{C+1}-T_C]
        =\E Q_C\,\E Q_C^{-1}\ge 1.
\end{equation}
Equality holds if and only if $p$ is uniform.
\end{proposition}

\begin{proof}
The identities are \eqref{eq:lru-hit-miss} and \eqref{eq:increment-mean}.  The inequality
is Cauchy's inequality, or Jensen's inequality applied to $q\mapsto q^{-1}$:
\[
        \E Q_C^{-1}\ge (\E Q_C)^{-1}.
\]
Equality holds if and only if $Q_C$ is almost surely constant.  Since $p_i>0$, every
$C$-subset occurs as $S_C$ with positive probability.  Thus $Q_C$ is constant if and only
if $p(A)$ is the same for every $C$-subset $A$.  Comparing two $C$-subsets that differ by
replacing $i$ with $j$ gives $p_i=p_j$ for all $i,j$.  Hence $p=u$.
\end{proof}

\section{The crossing method}

Fix a positive nonuniform vector $p\in\Delta_N^\circ$ and write
\[
        h_i=p_i-\frac1N,
        \qquad
        p_i(\theta)=\frac1N+\theta h_i,
        \qquad
        0\le\theta\le1.
\]
Then $\sum_i h_i=0$ and $p_i(\theta)>0$ for $0\le\theta\le1$.

Use a single family of independent standard Gumbel variables $G_1,\ldots,G_N$ for every
$\theta$, and define
\begin{equation}\label{eq:Ytheta}
        Y_i(\theta)=G_i+\log p_i(\theta).
\end{equation}
Let $S_C(\theta)$ be the set of indices corresponding to the $C$ largest values of
$Y_i(\theta)$.  By Lemma \ref{lem:gumbel}, $S_C(\theta)$ has the size-biased prefix law
with weights $p_\theta$.

Define the $h$-mass of the prefix by
\begin{equation}\label{eq:H-def}
        H_C(\theta)=\sum_{i\in S_C(\theta)}h_i.
\end{equation}
Then
\begin{equation}\label{eq:M-as-c-plus-theta-H}
        M_C(\theta)=p_\theta(S_C(\theta))
        =\frac CN+\theta H_C(\theta).
\end{equation}

\begin{lemma}[One-way crossings]\label{lem:one-way-crossings}
For $i\ne j$,
\begin{equation}\label{eq:gap-derivative}
        \frac{\dd}{\dd\theta}\bigl(Y_i(\theta)-Y_j(\theta)\bigr)
        =\frac{(1/N)(h_i-h_j)}{p_i(\theta)p_j(\theta)}.
\end{equation}
Consequently, if $h_i>h_j$, then item $i$ can overtake item $j$ as $\theta$ increases,
but item $j$ cannot overtake item $i$.
\end{lemma}

\begin{proof}
Since
\[
        \frac{\dd}{\dd\theta}\log p_i(\theta)=\frac{h_i}{p_i(\theta)},
\]
we have
\[
\begin{aligned}
        \frac{\dd}{\dd\theta}\bigl(Y_i-Y_j\bigr)
        &=\frac{h_i}{p_i(\theta)}-\frac{h_j}{p_j(\theta)}  \\
        &=\frac{h_ip_j(\theta)-h_jp_i(\theta)}{p_i(\theta)p_j(\theta)} \\
        &=\frac{(1/N)(h_i-h_j)}{p_i(\theta)p_j(\theta)}.
\end{aligned}
\]
The denominator is positive.  Thus the sign is the sign of $h_i-h_j$.
\end{proof}

\begin{lemma}[Pathwise monotonicity of prefix $h$-mass]\label{lem:H-pathwise}
For every fixed realization of the Gumbel variables outside a null set, the function
\[
        \theta\mapsto H_C(\theta)
\]
is nondecreasing on $[0,1]$.
\end{lemma}

\begin{proof}
Almost surely there are no triple ties and no pair is tied at a prescribed endpoint.  For
each pair $i,j$, Lemma \ref{lem:one-way-crossings} shows that their relative order can
change at most once, and when it changes the item with larger $h$ moves upward past the
item with smaller $h$.

The top-$C$ set changes only when a crossing occurs across the boundary between ranks $C$
and $C+1$.  At such a crossing, the entering item has larger $h$ than the leaving item.
Therefore the sum of $h_i$ over the top-$C$ set increases.  Crossings entirely inside the
top $C$ or entirely outside it leave the set, and hence $H_C$, unchanged.  Since there are
only finitely many pairs, $H_C(\theta)$ is a step function with nonnegative jumps.
\end{proof}

The next elementary lemma is the analytic step that converts pathwise monotonicity of
$H_C$ into increasing-convex order for $M_C$.

\begin{lemma}[A scaling lemma for increasing convex costs]\label{lem:scaling}
Let $X$ be a bounded real random variable with $\E X\ge0$.  Fix $c\in\R$ and an interval
$I$ containing $c+aX$ for all $a\in[a_0,a_1]$, where $0\le a_0\le a_1$.  If
$\phi:I\to\R$ is increasing and convex, then
\[
        a\mapsto \E\phi(c+aX)
\]
is nondecreasing on $[a_0,a_1]$.
\end{lemma}

\begin{proof}
On a compact interval, every increasing convex function admits the representation
\[
        \phi(x)=\alpha+\beta x+\int (x-t)_+\,\mu(\dd t),
\]
where $\beta\ge0$ and $\mu$ is a finite nonnegative measure.  It is therefore enough to
check the affine function $x$ and the hinge functions $(x-t)_+$.

For the affine function, monotonicity follows from $\E X\ge0$.  For a hinge, put
\[
        F_t(a)=\E(c+aX-t)_+.
\]
For $a>0$, the right derivative is
\[
        F_t'(a+)=\E\bigl[X\,\1_{\{X>(t-c)/a\}}\bigr].
\]
For any threshold $b$, the quantity $\E[X\1_{\{X>b\}}]$ is nonnegative.  If $b\ge0$ this
is immediate.  If $b<0$, then
\[
        \E[X\1_{\{X>b\}}]=\E X-\E[X\1_{\{X\le b\}}]\ge0,
\]
because $\E X\ge0$ and $X\1_{\{X\le b\}}\le0$.  Thus $F_t$ is nondecreasing.  The case
$a=0$ follows by right continuity.
\end{proof}

\section{The transform-order theorem}

\begin{theorem}[Radial increasing-convex order for cached mass]\label{thm:icx}
Let $N\ge2$, let $1\le C\le N-1$, and let $p\in\Delta_N^\circ$ be nonuniform.  Put
\[
        p_\theta=u+\theta(p-u),\qquad 0\le\theta\le1.
\]
Then for every increasing convex function $\phi$ on $[0,1]$, the map
\[
        \theta\mapsto \E\phi(M_C(p_\theta))
\]
is nondecreasing.  More explicitly, if $0\le\theta_1<\theta_2\le1$, then
\begin{equation}\label{eq:icx-order}
        \E\phi(M_C(p_{\theta_2}))
        \ge
        \E\phi(M_C(p_{\theta_1})).
\end{equation}
For every strictly increasing convex $\phi$, and also for the affine function
$\phi(x)=x$, the inequality is strict.
\end{theorem}

\begin{proof}
Use the common-Gumbel construction from Section 3.  Put
\[
        c=\frac CN,
        \qquad
        H_r=H_C(\theta_r),\quad r=1,2.
\]
By Lemma \ref{lem:H-pathwise}, $H_2\ge H_1$ pathwise.  Therefore, since $\phi$ is
increasing,
\begin{equation}\label{eq:first-step}
        \E\phi(c+\theta_2H_2)
        \ge
        \E\phi(c+\theta_2H_1).
\end{equation}

It remains to compare $\theta_2H_1$ and $\theta_1H_1$.  At $\theta=0$, the Gumbel order
is a uniformly random permutation, so $S_C(0)$ is a uniformly random $C$-subset.  Hence
\[
        \E H_C(0)=\frac CN\sum_i h_i=0.
\]
By Lemma \ref{lem:H-pathwise}, $H_1\ge H_C(0)$ pathwise, and therefore
\[
        \E H_1\ge0.
\]
Applying Lemma \ref{lem:scaling} with $X=H_1$, $a_0=\theta_1$, and $a_1=\theta_2$ gives
\begin{equation}\label{eq:second-step}
        \E\phi(c+\theta_2H_1)
        \ge
        \E\phi(c+\theta_1H_1).
\end{equation}
Combining \eqref{eq:first-step} and \eqref{eq:second-step}, and using
\eqref{eq:M-as-c-plus-theta-H}, proves \eqref{eq:icx-order}.

For strictness, choose $i,j$ with $h_i>h_j$.  Given any interval$(\theta_1,\theta_2)$, there is positive probability that $Y_i(\theta)-Y_j(\theta)$
crosses zero inside that interval.  Indeed the required condition is that
$G_i-G_j$ lies in a nonempty open interval, and $G_i-G_j$ has a positive density.  On a
further positive-probability event, exactly $C-1$ of the remaining Gumbels are so large
that they stay above both $i$ and $j$ throughout $[\theta_1,\theta_2]$, and all remaining
Gumbels are so small that they stay below both.  On this event the crossing changes the
membership of the top-$C$ set, replacing $j$ by $i$, and hence $H_C$ increases by
$h_i-h_j>0$.

Thus $\Pbb(H_2>H_1)>0$.  For strictly increasing $\phi$, the first comparison
\eqref{eq:first-step} is strict.  For $\phi(x)=x$, strictness follows from
\[
        \E M_C(\theta_2)-\E M_C(\theta_1)
        =\theta_2\E(H_2-H_1)+(\theta_2-\theta_1)\E H_1>0.
\]
This completes the proof.
\end{proof}

\begin{corollary}[Decreasing convex order for residual mass]\label{cor:dcx}
Under the hypotheses of Theorem \ref{thm:icx}, for every decreasing convex function
$\psi$ on $(0,1]$,
\begin{equation}\label{eq:decreasing-convex}
        \theta\mapsto \E\psi(Q_C(p_\theta))
\end{equation}
is nondecreasing whenever the expectation is finite.  It is strict for the strictly
decreasing costs considered below.
\end{corollary}

\begin{proof}
Put $\phi(m)=\psi(1-m)$.  If $\psi$ is decreasing and convex, then $\phi$ is increasing
and convex.  Since $Q_C=1-M_C$, Theorem \ref{thm:icx} gives the claim.
\end{proof}

\section{Consequences for LRU}

\subsection{Hit rate, miss rate, and stack-depth tails}

Taking $\phi(m)=m$ in Theorem \ref{thm:icx} gives the radial LRU theorem of
\cite{LongLRU2026}.

\begin{corollary}[LRU hit and miss rates]\label{cor:lru-hit-miss}
For $1\le C\le N-1$ and nonuniform $p\in\Delta_N^\circ$,
\[
        \theta\mapsto \Hit_{LRU}(C;p_\theta)=\E M_C(p_\theta)
\]
is strictly increasing on $[0,1]$, and
\[
        \theta\mapsto \MR_{LRU}(C;p_\theta)=\E Q_C(p_\theta)
\]
is strictly decreasing on $[0,1]$.
\end{corollary}

Equivalently, for the move-to-front search cost $R$ in stationarity,
\[
        \Pbb(R>C)=\MR_{LRU}(C;p).
\]
Therefore the stationary move-to-front search cost decreases in usual stochastic order
along every nonconstant ray away from uniform.  This is the radial restriction of the
Fill--Holst search-cost tail conjecture established in \cite{LongLRU2026}; the present
argument obtains it as the linear shadow of a stronger transform order.

\subsection{No-miss survival and hit-run lengths}

Taking $\phi(m)=m^n$ gives a stronger LRU statement.

\begin{corollary}[Discrete no-miss survival]\label{cor:no-miss-discrete}
For every integer $n\ge1$,
\[
        \theta\mapsto \E M_C(p_\theta)^n
\]
is strictly increasing.  Equivalently, starting from a stationary LRU cache of capacity
$C$, the probability that the next $n$ requests are all hits is strictly increasing along
every nonconstant ray away from uniform.
\end{corollary}

\begin{proof}
The function $m\mapsto m^n$ is increasing and convex on $[0,1]$.  Use
Theorem \ref{thm:icx} and identity \eqref{eq:no-miss-discrete}.
\end{proof}

Taking $\psi(q)=e^{-tq}$ gives the continuous-time version.

\begin{corollary}[Continuous no-miss survival]\label{cor:no-miss-continuous}
For every $t>0$,
\[
        \theta\mapsto \E e^{-tQ_C(p_\theta)}
\]
is strictly increasing.  Equivalently, in the rate-one Poissonized IRM, the probability
that no LRU miss occurs before time $t$ is strictly increasing along every nonconstant ray
away from uniform.
\end{corollary}

\begin{proof}
The function $q\mapsto e^{-tq}$ is decreasing and convex on $(0,1]$.  Apply
Corollary \ref{cor:dcx} and use \eqref{eq:no-miss-continuous}.
\end{proof}

\section{Consequences for coupon collecting}

\subsection{One-step stochastic monotonicity}

The coupon-collector increment $Y_C=T_{C+1}-T_C$ is conditionally geometric with success
probability $Q_C$.  Its survival probabilities are $\E M_C^n$.  Corollary
\ref{cor:no-miss-discrete} therefore gives an exact one-step stochastic order.

\begin{theorem}[Radial stochastic monotonicity of discovery increments]\label{thm:increment-stochastic}
For $0\le C<N$ and nonuniform $p\in\Delta_N^\circ$, the coupon-collector increment
\[
        T_{C+1}-T_C
\]
is stochastically increasing along the ray $p_\theta=u+\theta(p-u)$.  More precisely,
for every $n\ge0$,
\[
        \theta\mapsto \Pbb_{p_\theta}(T_{C+1}-T_C>n)
\]
is nondecreasing, and it is strict for $n\ge1$ when $1\le C\le N-1$.
\end{theorem}

\begin{proof}
For $n=0$ the survival probability is one.  For $n\ge1$, equation
\eqref{eq:geom-survival} gives
\[
        \Pbb(T_{C+1}-T_C>n)=\E M_C^n,
\]
and Corollary \ref{cor:no-miss-discrete} applies.  The case $C=0$ is deterministic with
$Q_0=1$ and $M_0=0$.
\end{proof}

In continuous time, the same statement follows from Corollary \ref{cor:no-miss-continuous}
and \eqref{eq:poissonized-increment-survival}.

\subsection{Reciprocal and negative moments}

\begin{corollary}[Negative residual moments and CCP increments]\label{cor:negative-moments}
For every $s>0$,
\[
        \theta\mapsto \E Q_C(p_\theta)^{-s}
\]
is strictly increasing for $1\le C\le N-1$.  In particular,
\begin{equation}\label{eq:reciprocal-increment-monotone}
        \theta\mapsto \E_{p_\theta}[T_{C+1}-T_C]
        =\E Q_C(p_\theta)^{-1}
\end{equation}
is strictly increasing.
\end{corollary}

\begin{proof}
For $s>0$, the function $q\mapsto q^{-s}$ is decreasing and convex on $(0,1]$.  Apply
Corollary \ref{cor:dcx}.  The identity \eqref{eq:reciprocal-increment-monotone} is
\eqref{eq:increment-mean}.
\end{proof}

Summing \eqref{eq:reciprocal-increment-monotone} over $C=0,\ldots,j-1$ recovers radial
monotonicity of the expected time to collect $j$ distinct coupons:
\[
        \theta\mapsto \E_{p_\theta}T_j
\]
is strictly increasing for $2\le j\le N$.  The theorem is stronger than this expectation
statement because it gives stochastic monotonicity of every one-step increment.

\section{Residual-mass cost classes}

Corollary \ref{cor:dcx} can be restated in cost language.  If $f$ is increasing and
concave on $[0,1]$, then $-f$ is decreasing and convex, so
\[
        \theta\mapsto \E f(Q_C(p_\theta))
\]
is nonincreasing.  If $g$ is completely monotone, then $g$ is decreasing and convex, so
\[
        \theta\mapsto \E g(Q_C(p_\theta))
\]
is nondecreasing.

\begin{corollary}[Fractional positive moments]\label{cor:fractional-positive}
For every $0<r\le1$,
\[
        \theta\mapsto \E Q_C(p_\theta)^r
\]
is strictly decreasing for $1\le C\le N-1$.
\end{corollary}

\begin{proof}
For $0<r<1$, the function $q\mapsto q^r$ is increasing and concave, equivalently
$q\mapsto -q^r$ is decreasing and convex.  The case $r=1$ is Corollary
\ref{cor:lru-hit-miss}.
\end{proof}

\begin{corollary}[Completely monotone residual costs]\label{cor:cm}
If $g$ is completely monotone on $(0,1]$ and finite on the range of $Q_C$, then
\[
        \theta\mapsto \E g(Q_C(p_\theta))
\]
is nondecreasing.  This includes $g(q)=e^{-\lambda q}$ and $g(q)=q^{-s}$.
\end{corollary}

\begin{proof}
A completely monotone function satisfies $g'\le0$ and $g''\ge0$, hence is decreasing and
convex.  Apply Corollary \ref{cor:dcx}.
\end{proof}

\subsection{Boundaries of the theorem}

The theorem deliberately does not assert monotonicity of all positive moments or centered
moments of $Q_C$.

\begin{example}[High positive moments can fail]\label{ex:high-positive}
For $C=1$,
\[
        Q_1=1-p_I,
        \qquad \Pbb(I=i)=p_i,
\]
and hence
\[
        \E Q_1^r=\sum_i p_i(1-p_i)^r.
\]
With $N=3$, $r=10$, and $p=(0.98,0.01,0.01)$,
\[
        \E_p Q_1^{10}
        =0.98(0.02)^{10}+0.02(0.99)^{10}
        \approx 0.0180876,
\]
whereas under uniform popularity
\[
        \E_u Q_1^{10}=(2/3)^{10}\approx0.0173415.
\]
Thus uniform does not maximize high positive residual moments.
\end{example}

\begin{example}[Centered residual moments are not radially monotone]\label{ex:centered}
Let $N=2$ and $C=1$.  Write
\[
        p_1=\frac12+\delta,
        \qquad
        p_2=\frac12-\delta.
\]
Then $Q_1$ equals $p_2$ with probability $p_1$ and $p_1$ with probability $p_2$.  A direct
calculation gives
\[
        \Var(Q_1)=\delta^2-4\delta^4.
\]
This increases near $\delta=0$ but decreases as $\delta$ approaches $1/2$.  Thus centered
moments of the residual mass are not governed by the radial transform order above.
\end{example}

\section{Berthet's relation revisited}

Berthet \cite{Berthet2017} observed that, for a partial coupon collection and an LRU
cache under the same IRM popularity, the miss rate and the forward difference of the
coupon waiting-time expectation are asymptotically related by
\[
        \MR_{LRU}(C;p)\,\bigl(\E T_{C+1}-\E T_C\bigr)\approx1.
\]
The exact residual identities show that the exact product is
\[
        \MR_{LRU}(C;p)\,\bigl(\E T_{C+1}-\E T_C\bigr)
        =\E Q_C\,\E Q_C^{-1}.
\]
Thus Berthet's relation is exact at uniform popularity, where $Q_C=(N-C)/N$ is
deterministic, and is approximately true whenever $Q_C$ is concentrated.

The transform order proved here gives a sharper finite-$N$ picture.  Along every
nonconstant ray away from uniform,
\[
        \E Q_C\downarrow,
        \qquad
        \E Q_C^{-1}\uparrow,
\]
and, more generally,
\[
        \E e^{-tQ_C}\uparrow,
        \qquad
        \E Q_C^{-s}\uparrow,
        \qquad
        \E Q_C^r\downarrow\quad(0<r\le1).
\]
LRU and coupon collecting are therefore linked not by a single approximate reciprocal
formula but by an exact common residual-hazard law.  LRU observes the residual mass
through the linear cost $q$; coupon collection observes the same residual mass through
the reciprocal cost $q^{-1}$.

\section{Relation to the positive pair-kernel LRU proof}

The radial LRU theorem in \cite{LongLRU2026} proves
\[
        \frac{\dd}{\dd\theta}\Hit_C(p_\theta)>0
\]
by deriving an explicit formula of the form
\[
        \frac{\dd}{\dd\theta}\Hit_C(p_\theta)
        =\frac{1}{N\theta}\sum_{a<b}(p_a(\theta)-p_b(\theta))^2
        K_{ab}^{(C)}(p_\theta),
        \qquad K_{ab}^{(C)}(p_\theta)>0.
\]
The present Gumbel proof gives a stronger conclusion but does not replace the value of
that local formula.  The two methods reveal different structures.

\begin{enumerate}[label=(\arabic*)]
\item The Gumbel crossing proof gives a global monotone coupling of the whole
Plackett--Luce ranking along a ray from uniform.  It proves increasing-convex order for
$M_C$ and gives the LRU--coupon transfer results.

\item The pair-kernel proof gives an explicit finite-dimensional derivative certificate
for the hit rate.  It is algebraic, computable, and useful for sensitivity analysis.
It also clarifies exactly why ray directions from uniform are special.

\item The occupancy-Jacobian proof in \cite{LongLRU2026} explains why full
Schur-convexity fails while radial monotonicity survives: the sign-indefinite local terms
collapse to positive squares only in the radial direction.
\end{enumerate}

Thus the new theorem subsumes the expectation-level conclusion of \cite{LongLRU2026} and
extends it to transform order, but the earlier proof remains valuable as a local analytic
and algebraic certificate.

\section{Further directions}

The results suggest several natural extensions.

\begin{enumerate}[label=(\arabic*)]
\item \textbf{Joint residual-mass order.}  The present theorem controls each fixed
$Q_C$ separately.  A stronger theorem would order the vector
\[
        (Q_0,Q_1,\ldots,Q_{N-1})
\]
or the vector of coupon-collector increments.  Such a result could imply stochastic
monotonicity of full partial collection times $T_j$, not merely their one-step increments.

\item \textbf{Other ranking paths.}  The crossing argument applies whenever scores have
the form $G_i+\eta_i(t)$ and the pairwise differences $\eta_i(t)-\eta_j(t)$ move
monotonically according to a fixed item score.  This includes softmax-temperature paths
$p_i(\beta)\propto e^{\beta v_i}$, and suggests a broader monotone-comparative-static
theory for Plackett--Luce rankings.

\item \textbf{Other caching policies.}  The argument is special to LRU/MTF because the
stationary order is Plackett--Luce.  Policies with different stationary cache laws, such
as CLIMB, would require different Gibbs, total-positivity, or polynomial methods.  The
positive-kernel and occupancy-Jacobian methods of \cite{LongLRU2026} may be more
adaptable there than the common-Gumbel crossing proof.

\item \textbf{Dispersion of the residual mass.}  The product
\[
        \E Q_C\,\E Q_C^{-1}
\]
is a Jensen dispersion gap.  The present theorem gives opposite monotonicities for the
two factors, but it does not assert monotonicity of the product.  Understanding this gap
more precisely would sharpen the interpretation of Berthet's approximation.
\end{enumerate}

\section*{Acknowledgment}
This paper was developed in response to the observation that the LRU miss rate and the
coupon-collector discovery increment are exact positive and reciprocal moment transforms
of the same size-biased residual mass.

\begin{thebibliography}{99}

\bibitem{AnceaumeBusnelSericola2015}
E. Anceaume, Y. Busnel and B. Sericola,
New results on a generalized coupon collector problem using Markov chains,
\emph{Journal of Applied Probability} \textbf{52} (2015), no. 2, 405--418.

\bibitem{AnceaumeBusnelSericola2016}
E. Anceaume, Y. Busnel, E. Schulte-Geers and B. Sericola,
Optimization results for a generalized coupon collector problem,
\emph{Journal of Applied Probability} \textbf{53} (2016), no. 2, 622--629.
Preprint: arXiv:1504.03878.

\bibitem{Berthet2017}
C. Berthet,
Contributions to the generalized coupon collector and LRU problems,
arXiv:1706.05250, 2017.

\bibitem{BonehHofri1997}
A. Boneh and M. Hofri,
The coupon-collector problem revisited: a survey of engineering problems and computational methods,
\emph{Stochastic Models} \textbf{13} (1997), no. 1, 39--66.

\bibitem{CheTungWang2002}
H. Che, Y. Tung and Z. Wang,
Hierarchical web caching systems: modeling, design and experimental results,
\emph{IEEE Journal on Selected Areas in Communications} \textbf{20} (2002), no. 7, 1305--1314.

\bibitem{Fagin1977}
R. Fagin,
Asymptotic miss ratios over independent references,
\emph{Journal of Computer and System Sciences} \textbf{14} (1977), no. 2, 222--250.

\bibitem{FaginPrice1978}
R. Fagin and T. G. Price,
Efficient calculation of expected miss ratios in the independent reference model,
\emph{SIAM Journal on Computing} \textbf{7} (1978), no. 3, 288--297.

\bibitem{FillHolst1996}
J. A. Fill and L. Holst,
On the distribution of search cost for the move-to-front rule,
\emph{Random Structures \& Algorithms} \textbf{8} (1996), no. 3, 179--186.

\bibitem{FlajoletGardyThimonier1992}
P. Flajolet, D. Gardy and L. Thimonier,
Birthday paradox, coupon collectors, caching algorithms and self-organizing search,
\emph{Discrete Applied Mathematics} \textbf{39} (1992), no. 3, 207--229.

\bibitem{Hildebrand1999}
M. Hildebrand,
On a conjecture of Fill and Holst involving the move-to-front rule and cache faults,
\emph{Probability in the Engineering and Informational Sciences} \textbf{13} (1999), no. 3, 377--385.

\bibitem{Jelenkovic1999}
P. R. Jelenkovic,
Asymptotic approximation of the move-to-front search cost distribution and least-recently-used caching fault probabilities,
\emph{The Annals of Applied Probability} \textbf{9} (1999), no. 2, 430--464.

\bibitem{King1971}
W. F. King III,
Analysis of demand paging algorithms,
\emph{Proceedings of IFIP Congress}, 1971, 485--490.

\bibitem{LongLRU2026}
C. D. Long,
Radial extremality for LRU caching and the Fill--Holst conjecture,
arXiv:2605.26107, 2026.

\bibitem{MarshallOlkinArnold2011}
A. W. Marshall, I. Olkin and B. C. Arnold,
\emph{Inequalities: Theory of Majorization and Its Applications},
2nd ed., Springer, New York, 2011.

\bibitem{VanichpunMakowski2004}
S. Vanichpun and A. M. Makowski,
Comparing strength of locality of reference: popularity, majorization, and some folk theorems,
\emph{Proceedings of IEEE INFOCOM 2004}, vol. 2, 838--849.

\end{thebibliography}

\end{document}